\section{守恒量引导的内禀尺度}
\subsection{守恒量的推导}
存在压力梯度的壁射流动量方程：
\begin{align}
    \frac{\partial uu}{\partial x}+\frac{\partial uv}{\partial y}=-\frac{1}{\rho}\frac{\partial p}{\partial x}+\nu\frac{\partial^{2} u}{\partial y^{2}}=-\frac{1}{\rho}\frac{\partial p}{\partial x}+\frac{1}{\rho}\frac{\partial\tau}{\partial y}
\end{align}
边界条件：
\begin{align}
    y=0,x>0&:\quad \psi=u=v=0\\
    y\rightarrow \infty，x>0&: \quad \tau=u=0
\end{align}

对动量方程从壁面法向某一点$y'$至$\infty$积分：
\begin{align}
    \int_{y'}^{\infty}\rho(\frac{\partial uu}{\partial x}+\frac{\partial uv}{\partial y})\mathrm{d}y'&=\int_{y'}^{\infty}\frac{\partial\tau}{\partial y}-\frac{\partial p}{\partial x}\mathrm{d}y'\\
    \frac{\partial}{\partial x}\rho\int_{y'}^{\infty}uu\mathrm{d}y'+\frac{\partial }{\partial x}\int_{y'}^{\infty}p\mathrm{d}y'-\rho uv&=-\tau
\end{align}

令$W=\int_{y'}^{\infty}(\rho uu+p)\mathrm{d}y'$，且有$\frac{\partial W}{\partial y}=-\rho uu-p$:
\begin{align}
    \frac{\partial W}{\partial x}-\rho uv+\tau=0
\end{align}
上式两边乘以$u$，与流场耦合
\begin{align}
     u\frac{\partial W}{\partial x}-\rho uuv+u\tau=0\\
     u\frac{\partial W}{\partial x}+v\frac{\partial W}{\partial y}+vp+u\tau=0
\end{align}
上面的式子描述了物理量$W$的对流输运。

将连续方程两边乘以$W$再加到上面的式子中：
\begin{align}
    W\frac{\partial u}{\partial x}+W\frac{\partial v}{\partial y}=0\\
    u\frac{\partial W}{\partial x}+v\frac{\partial W}{\partial y}+u\tau+vp+W\frac{\partial u}{\partial x}+W\frac{\partial v}{\partial y}=0\\
    \frac{\partial uW}{\partial x}+\frac{\partial vW}{\partial y}+vp+u\tau=0
\end{align}
如果补充非定常项:
\begin{align}
    \frac{\partial W}{\partial t}+\frac{\partial uW}{\partial x}+\frac{\partial vW}{\partial y}+vp+u\tau=0\\
    \frac{\mathrm{D}W}{\mathrm{D}t}+vp+u\tau=0
\end{align}
上式表明，广义机械能通量势（Generalized Mechanical Energy Potential）$W$ 的随体变化由两个机制驱动：一是主流速度 $u$与切应力 $\tau$ 的耦合作用，即壁面摩擦对流体的做功；二是压强通过法向速度 $v$ 所进行的功，即压强在法向上的作用，这项 $vp$ 描述了单位体积上由于压力梯度作用导致垂向流动所引入的压强功，是动量结构沿流向演化的重要推动因子。这个量描述了：壁射流中压力梯度+动量输运对局部能量结构的耦合控制机制

至此讨论的是在壁射流区域中某点$y$往上积分所得到的广义机械能通量势的输运方程

对上式在$y$方向上，从0至$\infty$积分：
\begin{align}
    \int_{0}^{\infty}\frac{\partial uW}{\partial x}\mathrm{d}y
    +\int_{0}^{\infty}\frac{\partial vW}{\partial y}\mathrm{d}y
    +\int_{0}^{\infty}(vp+u\tau)\mathrm{d}y=0
\end{align}
代入边界条件，最后得到：
\begin{align}
    \frac{\mathrm{d}}{\mathrm{d}x}\int_{0}^{\infty}uW\mathrm{d}y+\int_{0}^{\infty}vp\mathrm{d}y=0
\end{align}
假设：
\begin{align}
    \frac{\mathrm{d}Q}{\mathrm{d}x}=\int_{0}^{\infty}vp\mathrm{d}y
\end{align}
则有
\begin{align}
    \frac{\mathrm{d}}{\mathrm{d}x}\left(\int_{0}^{\infty}uW\mathrm{d}y+Q\right)=0
\end{align}
这表明：
\begin{align}
    \int_{0}^{\infty}uW\mathrm{d}y+Q=\int_{0}^{\infty}u\left(\int_{y'}^{\infty}(\rho uu+p)\mathrm{d}y'\right)\mathrm{d}y+Q=F=\cst
\end{align}
这说明在壁射流问题中$F$是一个守恒量。

\subsection{如何确定\(Q\)的具体表达式}

假设壁射流中最大流向速度所在的流线为势流，在此流线上引入Bernoulli方程重建压力梯度。

若最大速度所在的流线为 \(y = y_m(x)\)，从射流起点到该流线上任意一点有如下关系：
\begin{align}
    \frac{1}{2}u_m^2 + \frac{p(x, y_m)}{\rho} = \text{常数} \\
    \frac{\partial p(x, y_m)}{\partial x} = -\rho u_m \frac{\mathrm{d} u_m}{\mathrm{d} x}
\end{align}

假设：壁射流中压力的主要流向变化主要由射流中心（最大速度位置，记为 \(y = y_m(x)\)）上的压强分布主导，即
\begin{align}
    \frac{\partial p(x, y)}{\partial x} \approx \frac{\partial p(x, y_m)}{\partial x}
\end{align}

代入上一节的推导过程，有：
\begin{align}
    \frac{\mathrm{d}}{\mathrm{d}x} \int_0^\infty uW\,\mathrm{d}y - \int_0^\infty v\frac{u_m^2}{2}\,\mathrm{d}y = 0
\end{align}

记垂向体积流率：
\begin{align}
    V_i(x) = \int_0^\infty v\,\mathrm{d}y
\end{align}

于是：
\begin{align}
    \frac{\mathrm{d}Q}{\mathrm{d}x} = -\int_0^\infty v\frac{u_m^2}{2}\,\mathrm{d}y = -\frac{u_m^2(x)}{2} V_i(x)
\end{align}

这说明压力梯度项的贡献，可以以最大速度平方乘以垂向体积流率的形式体现。

该项 \( Q(x) \) 的物理意义在于：

\begin{itemize}
    \item 它表征了由于流向压力梯度的存在，使得射流流线在靠近壁面的区域发生法向弯折，从而形成非零的法向速度 \( v \)；
    \item 其中 \( u_m^2/2 \) 是射流中心处的动能密度，而 \( V_i(x) = \int_0^\infty v\,dy \) 描述了这一法向弯折趋势的强弱，反映了法向流动的整体量级；
    \item 因此，\( Q \sim \frac{u_m^2}{2} V_i \) 表示压力梯度引起的轴向动能通过法向速度发生“再分布”的程度，可视为压力梯度引导下流动结构的变化量。
\end{itemize}

该项对壁射流的流动特性具有重要影响：

\begin{enumerate}
    \item \textbf{调控自相似结构的扩展：}若 \( Q \) 项较大，说明法向偏折效应显著，主流动量需要通过更宽广的流动区域维持，从而导致射流更强烈地扩展；
    
    \item \textbf{改变最大速度 \( u_m(x) \) 的衰减速率：}\( Q \) 的存在抑制了 \( u_m \) 的快速衰减。在压力梯度作用下，经典的 \( u_m \sim x^{-a} \) 的幂律结构将受到改动，以维持整体动量守恒；
    
    \item \textbf{预测流动分离：}在逆压梯度（\( \mathrm{d}p/\mathrm{d}x > 0 \)）条件下，若 \( Q \) 增长速度超过动量扩散项增长，可能导致 \( \frac{\mathrm{d}}{\mathrm{d}x}F_{\text{eff}} < 0 \)，破坏自相似守恒结构，从而导致局部分离或回流；
    
    \item \textbf{扩展动量守恒形式：}传统动量积分中使用动量厚度项 \( \int uu\,dy \)，而在包含压力梯度的情形中，引入 \( Q \) 后的总动量守恒形式为：
    \[
        \frac{\mathrm{d}}{\mathrm{d}x} \left[ \int_0^\infty uW\,\mathrm{d}y + Q(x) \right] = 0
    \]
    这表明：总动量的守恒由剪切项与压力引导的动能转化项共同维护，为壁射流提供了更合理的物理描述。
\end{enumerate}

因此，\( Q \) 项在包含压力梯度的壁射流问题中不可忽略，它不仅提供了动能注入的来源，还决定了壁射流的扩展、速度衰减及是否产生分离的关键特征。



\subsection{壁射流问题中的内禀尺度}
一般的黏性流体问题中，至少存在两个守恒量，分别是$\rho$和$\mu$。根据上一小节的推导发现$F$是壁射流问题中特有的守恒量。现在我们使用这个三个守恒量重新构建壁射流的内禀尺度$\mathrm{U}$、$\mathrm{L}$，可以被用来确定壁射流相似解的具体形式。

壁射流守恒量的量纲表示为：
\begin{align}
    [F]&=\frac{\mathrm{ML^{2}}}{\mathrm{T}^{3}}\\
    [\rho]&=\frac{\mathrm{M}}{\mathrm{L^{3}}}\\
    [\mu]&=\frac{\mathrm{M}}{\mathrm{LT}}
\end{align}
上式中$\mathrm{M}$、$\mathrm{L}$、$\mathrm{T}$分别代表质量、长度、时间量纲。

现在用守恒量来表示这些基本量纲，然后用基本量纲来表示内禀尺度$\mathrm{U}$、$\mathrm{L}$。

基本量纲表示为：
\begin{align}
    \mathrm{M}&=\frac{\rho^{4}\nu^{9}}{F^{3}}\\
    \mathrm{L}&=\frac{\rho\nu^{3}}{F}\\
    \mathrm{T}&=\frac{\rho^{2}\nu^{5}}{F^{2}}
\end{align}

内禀尺度$U$、$L$的量纲分别为
\begin{align}
    [U]&=\frac{\mathrm{L}}{\mathrm{T}}\\
    [L]&=\mathrm{L}
\end{align}
用守恒量表示为：
\begin{align}
    U=\frac{F}{\rho\nu^{2}}\\
    L=\frac{\rho\nu^{3}}{F}
\end{align}

特征长度、特征速度、黏性系数组成的无量纲数表示为：
\begin{align}
    \frac{UL}{\nu}=1
\end{align}
代入以后发现满足这个无量纲数。